\chapter{Monte Carlo: simulare per capire}
\label{ch:metodo_montecarlo}

\epigraph{``In any conflict between two ideas, lapse into
random sampling and you'll learn surprising things.''}{adagio
del calcolo statistico}

\lettrine[lines=3]{D}{opo} aver prodotto un orario, una
domanda naturale del coordinatore \`e: \emph{quanto \`e
robusto rispetto a piccoli cambiamenti?} Cosa succede se un
docente cambia indisponibilit\`a all'ultimo minuto? Cosa
succede se una classe si sposta di sezione? Cosa succede se
viene aggiunta una compresenza dell'ultim'ora?

Una risposta secca \`e impossibile. Una risposta statistica,
invece, si pu\`o ottenere con il metodo \emph{Monte
Carlo}\index{Monte Carlo}: si \emph{perturba} l'input in
modo casuale, si \emph{rilancia} il solver molte volte, e
si \emph{guarda la distribuzione} dei risultati.

\section{La storia del nome}

\begin{storiabox}
Il metodo Monte Carlo nasce a Los Alamos negli anni '40
nel contesto del progetto Manhattan. Stanislaw Ulam, durante
una convalescenza, ide\`o di simulare l'evoluzione dei
neutroni in una bomba atomica con simulazioni stocastiche.
Lo chiam\`o ``Monte Carlo'' come scherzo, perch\'e suo zio
amava il casin\`o di Monte Carlo a Monaco. Metropolis e
Ulam pubblicarono il primo articolo formale nel
1949~\parencite{metropolis1949}.

L'idea si \`e poi diffusa in fisica, finanza, biologia,
statistica. Oggi \emph{Monte Carlo simulation} \`e un
metodo trasversale: ogni volta che un sistema deterministico
\`e troppo complesso per essere analizzato esattamente, si
campiona stocasticamente e si lavora sulle distribuzioni.
\end{storiabox}

\section{Sensitivity analysis: cosa misurare}

\begin{defbox}
La \emph{sensitivity
analysis}\index{sensitivity!analisi} consiste nel
\emph{perturbare} un parametro di un modello e nel misurare
quanto cambia l'output. Versione Monte Carlo: invece di
perturbare un parametro alla volta, perturba molti parametri
contemporaneamente con distribuzioni assegnate, e analizza
la distribuzione dell'output.
\end{defbox}

In $\pi$Tantum si offrono diverse sensitivit\`a, attivabili
dalla pagina Diagnostica:

\begin{itemize}
\item \textbf{Indisponibilit\`a docenti}: scegli $K$ scenari
  in cui un docente casuale acquisisce un'indisponibilit\`a
  in uno slot casuale. Per ognuno, rilancia la pipeline e
  registra: ha trovato soluzione? quanto \`e cambiato
  l'obiettivo? quante lezioni sono state spostate?
\item \textbf{Aggiunta compresenze}: simula l'arrivo
  dell'ultimo momento di una compresenza in una classe
  random per una materia random.
\item \textbf{Spostamento di una classe}: simula il cambio
  di curriculum/indirizzo di una classe.
\item \textbf{Generico}: l'utente passa una funzione di
  perturbazione personalizzata via DSL.
\end{itemize}

\section{Schema operativo}

\begin{enumerate}
\item Fissa un seed iniziale e un numero $K$ di campioni
  (default: 50).
\item Per ogni $i = 1, \ldots, K$:
  \begin{enumerate}
  \item Applica una perturbazione casuale $\delta_i$
    all'input.
  \item Rilancia la pipeline (Phase~A + B + C + D).
  \item Registra: feasibility, obiettivo finale, tempo,
    numero di lezioni cambiate rispetto al baseline.
  \end{enumerate}
\item Calcola statistiche descrittive: media, deviazione
  standard, mediana, percentili 5/95, frazione di scenari
  infeasible.
\item Visualizza la distribuzione (istogramma) del valore
  obiettivo.
\end{enumerate}

\subsection{Esempio numerico}

\begin{esempiobox}
Scuola \texttt{medium}: 28 classi, 50 docenti, baseline con
obiettivo $f_0 = 11.4$. Lanciamo Monte Carlo con $K = 50$
scenari di tipo ``indisponibilit\`a docente''.

Risultati: 47 scenari producono soluzione fattibile (3
infeasible); fra i 47, $\bar{f} = 12.1$, $\sigma = 0.8$,
mediana $= 12.0$, P5 $= 11.0$, P95 $= 13.6$. Il numero
medio di lezioni spostate \`e 6.3.

Lettura: il sistema \`e \emph{robusto in media} (l'obiettivo
peggiora di solo $\sim 6\%$), ma il 6\% degli scenari
diventa infeasible. Conviene investigare quali docenti sono
``critici'' (la cui indisponibilit\`a aggiuntiva rende il
problema infeasible).
\end{esempiobox}

\section{Quando \`e utile}

\begin{ideabox}
\begin{itemize}
\item \textbf{Pianificazione preventiva}: a settembre il
  coordinatore vuole sapere quali docenti sono ``critici''
  prima che si manifestino le assenze. Monte Carlo li
  identifica.
\item \textbf{Stima dell'impatto di una nuova richiesta}:
  un docente chiede di non lavorare il marted\`i. Monte
  Carlo simula l'aggiunta del vincolo (anche su altri
  docenti, per analogia) e dice quanto perde l'obiettivo
  in media.
\item \textbf{Validazione di una soluzione}: una soluzione
  con basso $f_0$ ma alta varianza nel Monte Carlo \`e
  \emph{fragile}; meglio cercarne una con $f_0$
  leggermente pi\`u alto ma varianza minore.
\end{itemize}
\end{ideabox}

\section{Quando non basta}

\begin{avvertenzabox}
\begin{itemize}
\item Monte Carlo \`e \emph{costoso}: $K$ esecuzioni della
  pipeline. Per scuole grandi $K = 50$ pu\`o richiedere
  ore. $\pi$Tantum lo lancia come run asincrono e mostra
  il progresso live.
\item L'analisi \`e \emph{indicativa}, non garantita: se
  $K$ \`e basso (es. 20), le statistiche hanno alta
  varianza. Aumentare $K$ migliora ma costa.
\item Monte Carlo \emph{non} dice perch\'e una perturbazione
  rompe il sistema; dice solo che lo rompe. Per capire
  perch\'e serve guardare i log dei singoli scenari
  (CP-SAT mostra il sottoinsieme di vincoli incompatibili).
\end{itemize}
\end{avvertenzabox}

\section{Connessioni}

\begin{connessionebox}
\begin{itemize}
\item \textbf{Hall pre-check} (cap.~\ref{ch:metodo_hall}):
  ogni scenario Monte Carlo passa prima per Hall; gli
  scenari che falliscono Hall sono contati come
  ``infeasible immediato'' senza dover lanciare CP-SAT.
\item \textbf{Statistica applicata}
  (cap.~\ref{ch:metodo_statistica}): le distribuzioni
  prodotte da Monte Carlo si analizzano con KS-test,
  istogrammi, percentili come qualsiasi altra
  distribuzione.
\end{itemize}
\end{connessionebox}

\begin{biblioparvabox}
\begin{itemize}
\item \cite{metropolis1949}: il primo paper formale.
  Cinque pagine.
\item \cite{robert2004}: \emph{Monte Carlo Statistical
  Methods}, manuale di riferimento moderno.
\item \cite{wasserman2004}: capp.~24--25 sull'inferenza
  Monte Carlo, accessibile.
\end{itemize}
\end{biblioparvabox}
